\documentclass{ximera}
% xmPreamble.tex en xmPrintstyle.sty worden automatisch geladen door ximera.cls (v2.7.8);
% NIET nog eens \input{./preamble.tex} of \addPrintStyle{.} doen (geeft 'already defined'-fouten).

\makeatletter % eigen img-map toevoegen aan het zoekpad (werkt standalone én vanuit de xourse)
\g@addto@macro\Ginput@path{{img/}{UW4202/img/}}
\makeatother

\begin{document}
\author{Alexander Holvoet}
\xmtitle{Zand op een niet-convexe veelhoek}{}

\begin{exercise}
    Teken de kamlijnen die je verwacht als je op deze figuur zand zou gieten.
    \begin{image}[7cm]
        \includegraphics[width=7cm,keepaspectratio]{zand26.jpg}
    \end{image}
    \begin{hint}
        Je eerste poging bestaat wellicht uit rechte lijnstukken (bissectrices). Verifieer
        met zand: in de binnenbocht klopt er iets niet. De zandhoop komt daar hoger dan je
        verwacht.
    \end{hint}
    \begin{oplossing}
        Veel leerlingen tekenen eerst enkel rechte lijnstukken.
        \begin{image}[12cm]
            \includegraphics[width=6cm,keepaspectratio]{zand27.jpg}
            \quad
            \includegraphics[width=6cm,keepaspectratio]{zand30.jpg}
        \end{image}
        Wie dit met zand verifieert, ziet dat de kamlijn niet alleen uit rechte lijnstukken
        bestaat: in zijaanzicht komt de zandhoop in de bocht wat hoger.
        \begin{image}[12cm]
            \includegraphics[width=5cm,keepaspectratio]{zand28.jpg}
            \quad
            \includegraphics[width=5cm,keepaspectratio]{zand29.jpg}
        \end{image}
        De stukken van de kamlijn kunnen niet allemaal rechtlijnig zijn. De punten tussen
        \(K\) en \(M\) (\(M\) zelf niet meegerekend) liggen dichter bij rand \(d\) dan bij
        punt \(F\). Het lijnstuk \([KM]\) moet vervangen worden door de punten die even ver
        van \(F\) als van rand \(d\) liggen: een stuk \emph{parabool} met brandpunt \(F\) en
        richtlijn \(d\). Analoog wordt \([KL]\) een stuk parabool met brandpunt \(F\) en
        richtlijn \(e\). Het lijnstuk \([KF]\) moet weggelaten worden: dit is geen kamlijn.
        Zandkorrels boven \([KF]\) twijfelen niet en rollen naar \(F\), het dichtste
        randpunt.
        \begin{image}[12cm]
            \includegraphics[width=6cm,keepaspectratio]{zand31.jpg}
            \quad
            \includegraphics[width=6cm,keepaspectratio]{zand32.jpg}
        \end{image}
        Met hoogtelijnen zie je dat de zandhoop bij het snijpunt van de twee paraboolstukken
        inderdaad iets hoger ligt dan de rechte kamlijnen.
    \end{oplossing}
\end{exercise}

\end{document}
